\chapter*{Como usar esta apostila}
\addcontentsline{toc}{chapter}{Como usar esta apostila}

Esta apostila foi escrita para quem está tendo o \textbf{primeiro contato com programação inteira}. Supomos
apenas que você já viu o básico de programação linear (variáveis, função objetivo, restrições, região viável) e que
sabe programar em Python. Tudo o mais é construído aqui, com calma e com exemplos.

O objetivo é que, ao final, você consiga resolver sozinho(a) todos os exercícios do notebook
\texttt{Aula\_Programacao\_Inteira\_Modelagem.ipynb} e, mais importante, \emph{entenda por que} cada modelo é
escrito do jeito que é.

\section*{Organização}

\begin{description}
  \item[Parte I -- Fundamentos.] O que é programação inteira, por que ela é necessária, por que ``resolver como PL e
  arredondar'' não funciona, o que é relaxação linear, como se formula um modelo, o que é o famoso \emph{big-M}, um
  pouco de lógica com variáveis binárias e um tutorial de PuLP (a biblioteca Python usada no notebook).
  \item[Parte II -- O notebook passo a passo.] Um capítulo por bloco do notebook (blocos 0 a 7). Em cada um:
  o problema, as respostas às perguntas da discussão, a construção do modelo matemático, a tradução para código, a
  interpretação do resultado e os erros mais comuns.
  \item[Apêndice.] Respostas dos exercícios de fixação e tabela com os valores ótimos esperados.
\end{description}

\section*{Como o notebook está organizado}

Cada bloco do notebook segue sempre a mesma dinâmica:
\begin{center}
\textbf{Discussão} $\;\rightarrow\;$ \textbf{Modelo} (no quadro) $\;\rightarrow\;$ \textbf{Exercício} (você completa os
\texttt{\# TODO}) $\;\rightarrow\;$ \textbf{Verificação}
\end{center}
As células de \textbf{Verificação} comparam o seu resultado com o gabarito \emph{sem revelar o valor}: aparece
``correto'' ou ``não bate com o gabarito'' (às vezes com uma dica). As células de \textbf{Solução de referência} ficam
ocultas; rodá-las sobrescreve a sua versão. Sugestão: só abra a solução depois de tentar, e use esta apostila
como ``professor particular'' quando travar.

\section*{Caixas usadas no texto}

\begin{definicao}[conceito formal]
Definições que você precisa saber de cor.
\end{definicao}
\begin{exemplo}[exemplo resolvido]
Pequenos problemas resolvidos do início ao fim, com todos os passos.
\end{exemplo}
\begin{receita}[técnica de modelagem]
Uma ``fórmula'' pronta para transformar uma regra em restrições lineares.
\end{receita}
\begin{atencao}[erros comuns]
Armadilhas em que quase todo mundo cai na primeira vez.
\end{atencao}
\begin{outros}[outras fontes]
Complementos de livros, artigos, cursos abertos e documentação de solvers, sempre com a referência
(lista completa no fim da apostila).
\end{outros}

\section*{Fontes}

A apostila segue o notebook da disciplina~\cite{notebook} e os slides \emph{Programação inteira --
Modelagem}~\cite{slides}, que por sua vez se baseiam em Wolsey~\cite{wolsey1998}, Arenales et
al.~\cite{arenales2007} e Sierksma~\cite{sierksma1996}. Para ampliar a visão, incluímos material de outros autores,
como o capítulo de programação inteira do livro do MIT de Bradley, Hax e Magnanti~\cite{amp}, artigos históricos
(Land e Doig, Gomory, Dantzig, Smith), surveys e a documentação de solvers comerciais. Todas as fontes estão nas
Referências.

\section*{Roteiro de estudo sugerido}

\begin{enumerate}
  \item Leia os Capítulos~1 a 4 antes da aula (ou antes de abrir o notebook).
  \item Faça os exercícios de fixação do fim de cada capítulo e confira no Apêndice.
  \item Leia o Capítulo~5 (PuLP) com o notebook aberto e rode as duas primeiras células.
  \item Para cada bloco do notebook: tente sozinho(a); se travar, leia a seção correspondente da Parte~II até o ponto em
  que destravar, e volte ao código.
  \item No fim, revise o quadro-resumo do Capítulo~14.
\end{enumerate}
